\input{preamble.tex}

\title{Лабораторная работа № 11. \\ Настройка NAT. Планирование}
\author{Сущенко Алина Николаевна \\ НПИбд-01-23}
\institute{Российский университет дружбы народов имени Патриса Лумумбы}
\date{2026}

\begin{document}

\frame{\titlepage}

\begin{frame}
\frametitle{Цель работы}
\begin{itemize}
    \item Провести подготовительные мероприятия по подключению локальной сети организации к Интернету.
    \item Построить схему соединения локальной сети с сетью провайдера и модельной сетью Интернета.
    \item Настроить IP-адреса и DNS-записи для последующей настройки NAT.
\end{itemize}
\end{frame}

\begin{frame}
\frametitle{Новая обновленная топология сети}
\begin{center}
    \includegraphics[width=0.7\textwidth]{../images/01.png}
\end{center}
\end{frame}

\begin{frame}
\frametitle{Создание новых зданий}
\begin{center}
    \includegraphics[width=0.7\textwidth]{../images/02.png}
\end{center}
\end{frame}

\begin{frame}
\frametitle{Оборудование в здании сети провайдера}
\begin{center}
    \includegraphics[width=0.9\textwidth]{../images/03.png}
\end{center}
\end{frame}

\begin{frame}
\frametitle{Оборудование в здании сети модельного Интернета}
\begin{center}
    \includegraphics[width=0.9\textwidth]{../images/04.png}
\end{center}
\end{frame}

\end{frame}

\begin{frame}
\frametitle{Настройка модулей медиаконвертера}
\begin{center}
    \includegraphics[width=0.9\textwidth]{../images/05.png}
\end{center}

\end{frame}

\begin{frame}[fragile]
\frametitle{Настройка маршрутизатора ansusenko-gw-1}

\begin{center}
    \includegraphics[width=0.7\textwidth]{../images/06.png}
\end{center}
\end{frame}

\begin{frame}
\frametitle{Настройка IP-адреса сервера www.yandex.ru}
\begin{center}
    \includegraphics[width=0.9\textwidth]{../images/07.png}
\end{center}
\end{frame}

\begin{frame}
\frametitle{Настройка IP-адреса сервера stud.rudn.university}
\begin{center}
    \includegraphics[width=0.9\textwidth]{../images/08.png}
\end{center}
\end{frame}

\begin{frame}
\frametitle{Настройка IP-адреса сервера esystem.pfur.ru}
\begin{center}
    \includegraphics[width=0.9\textwidth]{../images/09.png}
\end{center}
\end{frame}

\begin{frame}
\frametitle{Настройка IP-адреса сервера www.rudn.ru}
\begin{center}
    \includegraphics[width=0.9\textwidth]{../images/10.png}
\end{center}
\end{frame}

\begin{frame}
\frametitle{Настройка DNS-сервера}
\begin{center}
    \includegraphics[width=0.95\textwidth]{../images/11.png}
\end{center}
\end{frame}

\begin{frame}
\frametitle{Схема L1 - таблица портов подключения}
\begin{center}
\small
\begin{tabular}{|l|l|l|l|}
\hline
\textbf{Устройство 1} & \textbf{Порт} & \textbf{Устройство 2} & \textbf{Порт} \\
\hline
msk-ansusenko-gw-1 & fa0/1 & ansusenko-mc-2 & fa0 \\
ansusenko-mc-2 & fa1 & ansusenko-mc-1 & fa1
\\
ansusenko-mc-1 & fa0 & ansusenko-sw-1 & fa0/2 \\
ansusenko-sw-1 & fa0/1 & ansusenko-gw-1 & fa0/0 \\
ansusenko-gw-1 & fa0/1 & ansusenko-mc-02 & fa0 \\
ansusenko-mc-02 & fa1 & ansusenko-mc-01 & fa1 \\
ansusenko-mc-01 & fa0 & internet-sw-1 & fa0/1 \\
internet-sw-1 & fa0/2 & esystem.pfur.ru & fa0 \\
internet-sw-1 & fa0/3 & www.rudn.ru & fa0 \\
internet-sw-1 & fa0/4 & stud.rudn.university & fa0 \\
internet-sw-1 & fa0/5 & www.yandex.ru & fa0 \\
\hline
\end{tabular}
\end{center}
\end{frame}

\begin{frame}
\frametitle{Схема L3 - таблица IP-адресов}
\begin{center}
\small
\begin{tabular}{|l|l|l|l|}
\hline
\textbf{Устройство} & \textbf{Интерфейс} & \textbf{IP-адрес} & \textbf{Маска} \\
\hline
ansusenko-gw-1 & fa0/0 & 10.128.0.1 & 255.255.255.0 \\
ansusenko-gw-1 & fa0/1 & 192.0.2.1 & 255.255.255.0 \\
www.yandex.ru & fa0 & 192.0.2.11 & 255.255.255.0 \\
stud.rudn.university & fa0 & 192.0.2.12 & 255.255.255.0 \\
esystem.pfur.ru & fa0 & 192.0.2.13 & 255.255.255.0 \\
www.rudn.ru & fa0 & 192.0.2.14 & 255.255.255.0 \\
ansusenko-dns & — & 10.128.0.5 & 255.255.255.0 \\
\hline
\end{tabular}
\end{center}
\end{frame}

\begin{frame}
\frametitle{Выводы}
\begin{itemize}
    \item В ходе выполнения лабораторной работы была построена физическая топология L1, включающая локальную сеть организации, сеть провайдера (медиаконвертеры \texttt{ansusenko-mc-1}, \texttt{ansusenko-mc-2}) и модельную сеть Интернета (коммутатор \texttt{internet-sw-1}, 4 сервера).Настроены IP-адреса на маршрутизаторе \texttt{ansusenko-gw-1}. Подготовлена инфраструктура для последующей настройки механизма NAT.
\end{itemize}
\end{frame}

\begin{frame}
\frametitle{Контрольные вопросы}
\begin{block}{Вопрос 1. Что такое Network Address Translation (NAT)?}
NAT (Network Address Translation) - механизм преобразования IP-адресов транзитных пакетов. Используется для обеспечения доступа устройств локальных сетей с внутренними (частными) IP-адресами к сети Интернет путём замены частного адреса на публичный.
\end{block}
\end{frame}

\begin{frame}
\frametitle{Контрольные вопросы}
\begin{block}{Вопрос 2. Как определить, находится ли узел сети за NAT?}
Определить, находится ли узел за NAT, можно следующими способами:
\begin{itemize}
    \item Сравнить IP-адрес, который видит удалённый сервер, с IP-адресом, указанным на локальном интерфейсе узла — если они различаются, узел за NAT.
    \item Использовать протоколы, которые не проходят через NAT (например, некоторые версии SIP, FTP в активном режиме).
    \item Проанализировать поле TTL в IP-пакетах - при прохождении NAT TTL обычно не меняется, но косвенные методы могут указать на наличие NAT.
\end{itemize}
\end{block}
\end{frame}

\begin{frame}
\frametitle{Контрольные вопросы}
\begin{block}{Вопрос 3. Какое оборудование отвечает за преобразование адреса методом NAT?}
Преобразование адресов методом NAT выполняют:
\begin{itemize}
    \item Маршрутизаторы (например, Cisco 2811 в данной работе).
    \item Межсетевые экраны (firewalls).
    \item Специализированные устройства (NAT-шлюзы, роутеры провайдеров).
    \item Операционные системы серверов (например, iptables в Linux).
\end{itemize}
В данной работе NAT настраивается на маршрутизаторе \texttt{ansusenko-gw-1}.
\end{block}
\end{frame}

\begin{frame}
\frametitle{Контрольные вопросы}
\begin{block}{Вопрос 4. В чём отличие статического, динамического и перегруженного NAT?}
\begin{itemize}
    \item \textbf{Статический NAT (SNAT)} - преобразование 1:1. Каждому внутреннему адресу постоянно соответствует один внешний адрес.
    \item \textbf{Динамический NAT (DNAT)} - преобразование 1:N. Внутренний адрес получает один внешний из пула доступных адресов (но не более одного внешнего на один внутренний).
    \item \textbf{Перегруженный NAT (NAT Overload / PAT)} - преобразование N:1. Множество внутренних адресов используют один внешний адрес, различаясь номерами портов (Port Address Translation).
\end{itemize}
\end{block}
\end{frame}

\begin{frame}
\frametitle{Контрольные вопросы}
\begin{block}{Вопрос 5. Охарактеризуйте типы NAT.}
Основные типы NAT:
\begin{itemize}
    \item \textbf{Static NAT (SNAT)} - постоянное отображение внутреннего адреса на внешний. Используется для серверов, доступных извне.
    \item \textbf{Dynamic NAT} - временное отображение из пула внешних адресов. Требует достаточно внешних адресов.
    \item \textbf{PAT (Port Address Translation)} - самый распространённый тип. Использует один внешний адрес для всей локальной сети, различая соединения по портам.
    \item \textbf{Double NAT} - последовательное применение NAT (например, у провайдера и в локальной сети).
\end{itemize}
В данной работе планируется использование PAT (NAT Overload) на маршрутизаторе \texttt{ansusenko-gw-1}.
\end{block}
\end{frame}

\end{document}
